\documentclass[12pt]{article}
\usepackage[margin=0.58in,top=0.62in,bottom=0.48in]{geometry}
\usepackage{lmodern}
\usepackage{microtype}
\usepackage{fancyhdr}
\usepackage{titlesec}
\usepackage{enumitem}
\usepackage{lastpage}
\usepackage[hidelinks]{hyperref}

\setlength{\parindent}{0pt}
\setlength{\parskip}{0.18em}
\setlength{\headheight}{26pt}
\raggedbottom
\titleformat{\section}{\normalsize\bfseries\scshape}{}{0pt}{}
\titleformat{\subsection}{\normalsize\bfseries}{}{0pt}{}
\titlespacing*{\section}{0pt}{0.35em}{0.12em}
\titlespacing*{\subsection}{0pt}{0.25em}{0.08em}
\pagestyle{fancy}
\fancyhf{}
\fancyhead[C]{\footnotesize Lit Example Reader \quad Assignment ID: U1L11LE \quad Page \thepage\ of \pageref{LastPage}}
\renewcommand{\headrulewidth}{0.5pt}

\newenvironment{playpassage}{\begin{quote}\footnotesize\setlength{\parskip}{0.08em}\setlength{\parindent}{0pt}}{\end{quote}}
\newcommand{\speaker}[1]{\textbf{#1.}}

\begin{document}

\begin{center}
{\LARGE\bfseries Week 11 Lit Example Reader}\par\vspace{0.02em}
{\large\scshape Julius Caesar and Macbeth: Testing Comparisons}\par\vspace{0.06em}
\rule{0.78\textwidth}{0.5pt}
\end{center}

\textbf{Purpose:} Shakespeare passages for Week 11: test whether a comparison proves the claim or only makes it vivid.

\textbf{What to watch for:} Source case, target case, relevant similarity, important disanalogy, and overextension.

\section*{Before the Passages: The Week 11 Logic Tool}

\begin{itemize}[leftmargin=1.3em,itemsep=0.05em,topsep=0.08em,parsep=0.03em]
\item \textbf{Source case:} the comparison used to illuminate the argument.
\item \textbf{Target case:} the person, choice, or situation being judged.
\item \textbf{Relevant similarity:} the likeness that matters for the conclusion.
\item \textbf{Disanalogy:} an important difference that limits the proof.
\end{itemize}

\section*{Passage 1: Julius Caesar, Act 2, Scene 1 --- Brutus's Serpent's Egg}

\textbf{Context:} Brutus is deciding whether Caesar should be killed before Caesar has actually become king. Watch the comparison he uses to reason about future danger.

\begin{playpassage}
\speaker{Brutus} It must be by his death: and for my part,\\
I know no personal cause to spurn at him,\\
But for the general. He would be crown'd:\\
How that might change his nature, there's the question.\\
It is the bright day that brings forth the adder;\\
And that craves wary walking. Crown him? That;\\
And then, I grant, we put a sting in him,\\
That at his will he may do danger with.\\
The abuse of greatness is, when it disjoins\\
Remorse from power: and, to speak truth of Caesar,\\
I have not known when his affections sway'd\\
More than his reason. But 'tis a common proof,\\
That lowliness is young ambition's ladder,\\
Whereto the climber-upward turns his face;\\
But when he once attains the upmost round,\\
He then unto the ladder turns his back.\\
...\\
And therefore think him as a serpent's egg\\
Which hatch'd, would, as his kind, grow mischievous,\\
And kill him in the shell.
\end{playpassage}

\subsection*{Reader Notes}

This is your assisted example. Brutus's target case is Caesar with crown-like power. His source case is a serpent's egg: not dangerous yet, but dangerous by nature if allowed to hatch. The relevant similarity Brutus wants is future danger before visible harm. The disanalogy is that Caesar is a human political actor, not an animal whose nature is fixed and known. Week 11 asks whether the analogy proves assassination is necessary or only supports caution about power.

\textbf{Reading task:} Underline the serpent's egg sentence. Circle one line where Brutus admits uncertainty about Caesar. In the margin, write the conclusion Brutus wants the comparison to support.

\newpage

\section*{Passage 2: Macbeth, Act 5, Scene 7 --- Macbeth as Bear at Bay}

\textbf{Context:} Near the end of the play, Macbeth compares himself to a trapped animal. Watch what the comparison explains, and what it does not prove.

\begin{playpassage}
\speaker{Macbeth} They have tied me to a stake; I cannot fly,\\
But, bear-like, I must fight the course.
\end{playpassage}

\subsection*{Reader Notes}

This passage shifts more work to you. Macbeth's comparison can explain pressure, desperation, and forced combat. Mark the source case and target case. Then decide whether the image proves that Macbeth has no moral choice left, or whether it mainly describes how trapped he feels.

\textbf{Reading task:} Label the source case and target case. Next to the passage, write one relevant similarity and one possible disanalogy. Save the final analogy judgment for the worksheet.

\section*{How to Use These Passages on the Worksheet}

Do not summarize the plots. Use Week 11's test: what comparison is doing the proof work, and does the relevant similarity actually carry the conclusion?

\end{document}
